\section{Deterministic finite automata}

\subsection{Formal definition}
\begin{definition}[DFA]
A \emph{deterministic finite automaton} is a $5$-tuple
\[
  M=(Q,\Sigma,\transition,q_0,F)
\]
where
\begin{enumerate}
  \item $Q$ is a finite set of states;
  \item $\Sigma$ is a finite alphabet;
  \item $\transition:Q\times\Sigma\to Q$ is the \emph{transition function};
  \item $q_0\in Q$ is the start state;
  \item $F\subseteq Q$ is the set of accept states.
\end{enumerate}
\end{definition}

Two remarks:
\begin{enumerate}
  \item $F$ may be empty. Then $\lang{M}=\varnothing$: no string is accepted.
  \item Determinism means that for every pair $(q,a)\in Q\times\Sigma$ there is
        \emph{exactly one} next state $\transition(q,a)$.
\end{enumerate}

\begin{example}[The $5$-tuple for $M_1$]
\[
\begin{array}{rl}
  Q&=\set{q_1,q_2,q_3},\\
  \Sigma&=\set{0,1},\\
  q_0&=q_1,\\
  F&=\set{q_2},
\end{array}
\qquad
\begin{array}{c|cc}
\transition & 0 & 1\\
\hline
q_1 & q_1 & q_2\\
q_2 & q_3 & q_2\\
q_3 & q_2 & q_2
\end{array}
\]
If $M$ accepts a set $A$, we call $A$ the \emph{language of} $M$ and write
$A=\lang{M}$.
\end{example}

\subsection{A modulo-3 machine with RESET}
\begin{example}[Running sum modulo $3$]
Let $\Sigma=\set{0,1,2,\mathtt{RESET}}$. Think of $0,1,2$ as adding that
value to a running total, and of $\mathtt{RESET}$ as setting the total back
to $0$. Accept if the current total is $0\bmod 3$.

States $q_0,q_1,q_2$ store the current residue. The start and accept state is
$q_0$. Transitions on $1$ and $2$ move around the triangle; $0$ is a self-loop;
$\mathtt{RESET}$ (written $\mathtt{R}$) returns to $q_0$ from every state.
\begin{center}
\begin{tikzpicture}[automaton]
  \node[state,initial,accepting] (q0) at (0,0) {$q_0$};
  \node[state] (q1) at (3.4,1.7) {$q_1$};
  \node[state] (q2) at (3.4,-1.7) {$q_2$};
  \path[->]
    (q0) edge[loop above] node{$0$} ()
         edge[bend left=12] node[above]{$1$} (q1)
         edge[bend right=12] node[below]{$2$} (q2)
    (q1) edge[loop above] node{$0$} ()
         edge[bend left=12] node[right]{$1$} (q2)
         edge[bend left=18] node[above]{$2$} (q0)
    (q2) edge[loop below] node{$0$} ()
         edge[bend left=18] node[below]{$1$} (q0)
         edge[bend left=12] node[right]{$2$} (q1)
    (q1) edge[bend right=32] node[below left]{\texttt{R}} (q0)
    (q2) edge[bend left=32] node[above left]{\texttt{R}} (q0)
    (q0) edge[loop below] node{\texttt{R}} ();
\end{tikzpicture}
\end{center}
Finite memory is enough because only the residue class of the sum matters.
\end{example}

\subsection{Designing finite automata}
Design principle:
\begin{itemize}
  \item You must be ready to accept or reject after every prefix.
  \item You do \emph{not} have infinite memory.
  \item Each state records exactly the information you still need about the
        prefix read so far.
\end{itemize}

\begin{example}[Odd number of $1$s]
$A=\set{w\in\set{0,1}^*:w\text{ has an odd number of }1\text{s}}$.
Remember only the parity of the number of $1$s.
\begin{center}
\begin{tikzpicture}[automaton,node distance=34mm]
  \node[state,initial] (even) {even};
  \node[state,accepting,right=of even] (odd) {odd};
  \path[->]
    (even) edge[loop above] node{$0$} ()
          edge[bend left=16] node[above]{$1$} (odd)
    (odd) edge[loop above] node{$0$} ()
         edge[bend left=16] node[below]{$1$} (even);
\end{tikzpicture}
\end{center}
\end{example}

\begin{example}[Substring $001$]
$A=\set{w\in\set{0,1}^*:w\text{ contains }001\text{ as a substring}}$.
States record the longest prefix of $001$ that is a suffix of the text so
far:
\[
  q_0:\text{nothing},\quad
  q_1:\text{suffix }0,\quad
  q_2:\text{suffix }00,\quad
  q_3:\text{already seen }001.
\]
\begin{center}
\begin{tikzpicture}[automaton,node distance=26mm]
  \node[state,initial] (q0) {$q_0$};
  \node[state,right=of q0] (q1) {$q_1$};
  \node[state,right=of q1] (q2) {$q_2$};
  \node[state,accepting,right=of q2] (q3) {$q_3$};
  \path[->]
    (q0) edge[loop above] node{$1$} ()
         edge node[above]{$0$} (q1)
    (q1) edge[bend left=18] node[above]{$1$} (q0)
         edge node[above]{$0$} (q2)
    (q2) edge[loop above] node{$0$} ()
         edge node[above]{$1$} (q3)
    (q3) edge[loop above] node{$0,1$} ();
\end{tikzpicture}
\end{center}
On a $0$ from $q_2$ we stay in $q_2$, because the suffix is still $00$.
\end{example}

\subsection{Formal definition of computation}
\begin{definition}[Computation of a DFA]
Let $M=(Q,\Sigma,\transition,q_0,F)$ and let $w=w_1w_2\cdots w_n$ with each
$w_i\in\Sigma$ (the case $n=0$ is $w=\eps$). Then $M$ \emph{accepts} $w$ if
there is a sequence of states $r_0,r_1,\ldots,r_n\in Q$ such that
\begin{enumerate}
  \item $r_0=q_0$;
  \item $\transition(r_i,w_{i+1})=r_{i+1}$ for $0\leq i\leq n-1$;
  \item $r_n\in F$.
\end{enumerate}
\end{definition}
Equivalently, extend $\transition$ to strings by
\[
  \transition^*(q,\eps)=q,\qquad
  \transition^*(q,wa)=\transition(\transition^*(q,w),a).
\]
Then $\lang{M}=\set{w\in\Sigma^*:\transition^*(q_0,w)\in F}$.

If $\transition(r_i,w_{i+1})=r_{i+1}$, we say that $M$ is in state $r_{i+1}$
after processing $w_1\cdots w_{i+1}$.

\begin{definition}[Regular language]
A language is \emph{regular} if some DFA recognizes it.
\end{definition}
